\GalSourcePageBoundary{074}{L0073}{45}{45}

voici l'un de ces groupes:
\[
\begin{array}{ccc}
  abcd, & acdb, & adbc,  \\
  badc, & cabd, & dacb,  \\
  cdab, & dbac, & bcad,  \\
  dcba, & bdca, & cbda.
\end{array}
\]
Maintenant ce groupe se partage lui-même en trois groupes,
comme il est indiqué aux théorèmes II et III. Ainsi, par l'extraction
d'un seul radical du troisième degré, il reste simplement le
groupe
\[
\begin{array}{l}
  abcd,  \\
  badc,  \\
  cdab,  \\
  dcba:
\end{array}
\]
ce groupe se partage de nouveau en deux groupes:
\[
\begin{array}{ll}
  abcd, & cdab,  \\
  badc, & dcba.
\end{array}
\]
Ainsi, après une simple extraction de racine carrée, il restera
\[
\begin{array}{l}
  abcd,  \\
  badc;
\end{array}
\]
ce qui se résoudra enfin par une simple extraction de racine
carrée.

On obtient ainsi, soit la solution de Descartes, soit celle
d'Euler; car, bien qu'après la résolution de l'équation auxiliaire
du troisième degré ce dernier extraye trois racines carrées, on
sait qu'il suffit de deux, puisque la troisième s'en déduit rationnellement.

Nous allons maintenant appliquer cette condition aux équations
irréductibles dont le degré est premier.
